%!TEX TS-program = xelatex
%!TEX encoding = UTF-8 Unicode
\documentclass[10pt, a4paper]{article}

\usepackage[top=0.8cm, bottom=0.8cm, left=1.2cm, right=1.2cm]{geometry}
\usepackage{fontspec}
\usepackage[default,opentype]{sourcesanspro}
\usepackage{fontawesome}
\usepackage{hyperref}
\usepackage{xcolor}
\usepackage{titlesec}
\usepackage{enumitem}
\usepackage{tabularx}
\usepackage{fancyhdr}

% Colors
\definecolor{headercolor}{HTML}{000000}
\definecolor{subheadercolor}{HTML}{000000}
\definecolor{datecolor}{HTML}{000000}

% Hyperlink setup
\hypersetup{colorlinks=true, urlcolor=black, linkcolor=black}

% Auto month-year of last compilation (updates on every CI build)
\newcommand{\monthyear}{\ifcase\month\or January\or February\or March\or April\or May\or June\or July\or August\or September\or October\or November\or December\fi\space\the\year}

% Remove page numbers / set footer
\pagestyle{fancy}
\fancyhf{}
\renewcommand{\headrulewidth}{0pt}
\fancyfoot[C]{\footnotesize\textcolor{datecolor}{Fabio Cavagna~~~\textperiodcentered~~~Curriculum Vitae~~~\textperiodcentered~~~Last updated: \monthyear}}

% Section formatting
\titleformat{\section}{\large\bfseries\color{headercolor}\uppercase}{}{0em}{}[\titlerule]
\titlespacing*{\section}{0pt}{0.4em}{0.3em}

% Custom commands
\newcommand{\cventry}[5]{%
  \noindent\begin{tabularx}{\textwidth}{@{}X r@{}}
    \textbf{#2} & #3 \\
    \textit{#1} & \textcolor{datecolor}{#4} \\
  \end{tabularx}
  \ifx&#5&\else
    \vspace{-0.5em}
    #5
  \fi
  \vspace{0.1em}
}

\newcommand{\cvskillrow}[2]{%
  \textbf{#1:} & #2 \\[0.2em]
}

\setlist[itemize]{leftmargin=1.5em, itemsep=0pt, parsep=0pt, topsep=1pt}

\begin{document}

%-------------------------------------------------------------------------------
% HEADER
%-------------------------------------------------------------------------------
\begin{center}
  {\LARGE\bfseries Fabio Cavagna}\\[0.3em]
  {\textit{PhD Candidate in Aerospace Engineering}}\\[0.4em]
  \faEnvelope~\href{mailto:fabiocavagna0@gmail.com}{fabiocavagna0@gmail.com} \quad
  \faLinkedin~\href{https://www.linkedin.com/in/fabio-cavagna}{fabio-cavagna}
\end{center}

\vspace{-0.8em}

%-------------------------------------------------------------------------------
% EXPERIENCE
%-------------------------------------------------------------------------------
\section{Experience}

\cventry
  {Research Test and Simulation Engineer}
  {Siemens Digital Industries Software}
  {Leuven, Belgium}
  {Dec. 2024 -- Present}
  {\begin{itemize}
    \item Conducted experimental tests of structures mounted on shaker tables, post-processing the acquired data to extract modal parameters with Simcenter Testlab.
    \item Performed numerical simulations of structures using Simcenter Nastran, validating the results against experimental data.
  \end{itemize}}

\cventry
  {Structural Dynamics Intern}
  {CNES}
  {Paris, France}
  {Apr. 2024 -- Oct. 2024}
  {\begin{itemize}
    \item Defined testing procedures for the experimental campaign, including the finalization of the installation of a small-scale tank and the definition of safety measures.
    \item Performed numerical simulations to compute the sloshing and hydro-elastic modes of liquid tanks.
  \end{itemize}}

\cventry
  {Master Thesis Intern}
  {Thales Alenia Space}
  {Turin, Italy}
  {Mar. 2023 -- Sept. 2023}
  {\begin{itemize}
    \item Conducted a comprehensive assessment of methodologies used to evaluate and compare the severity of mechanical loads, including Shock Response Spectrum, Extreme Response Spectrum, and Fatigue Damage Spectrum.
    \item Developed a novel approach for comparing the severity of mechanical Multi Degree of Freedom loads using MATLAB, emphasizing the analysis of interface forces in the presence of 6 DoF loads.
  \end{itemize}}

\cventry
  {Structural Engineer}
  {Polispace}
  {Milan, Italy}
  {Mar. 2023 -- Sept. 2023}
  {\begin{itemize}
    \item Student association developing a CubeSat for ESA's ``Fly Your Satellite!'' program. Simplified CAD geometries using SolidWorks and Abaqus, integrating them into the FEM model.
    \item Contributed to defining structural requirements for the Thermal Vacuum Chamber Test to ensure CubeSat resilience in extreme conditions.
  \end{itemize}}

%-------------------------------------------------------------------------------
% EDUCATION
%-------------------------------------------------------------------------------
\section{Education}

\cventry
  {PhD Candidate in Aerospace Engineering}
  {Politecnico di Milano}
  {Milan, Italy}
  {Dec. 2024 -- Present}
  {Research topic: ``Decoupling techniques based on substructuring for fixed-base analysis.''}

\vspace{0.6em}

\cventry
  {Advanced Master in Aeronautical and Space Structures}
  {ISAE-SUPAERO}
  {Toulouse, France}
  {Oct. 2023 -- Oct. 2024}
  {\begin{itemize}
    \item Team Innovative Management for Evolved Strategies -- TRL analysis for LH2 propellant transitioning from space to aeronautics. Defined WBS, Gantt chart, and risk matrix.
    \item Space Structures -- Sizing of the TMI Wall of SPOT 4 for quasi-static loads, thermal loads, and buckling through Patran/MSC Nastran.
  \end{itemize}}

\vspace{0.6em}

\cventry
  {M.S. in Space Engineering}
  {Politecnico di Milano}
  {Milan, Italy}
  {Sept. 2021 -- Oct. 2023}
  {\begin{itemize}
    \item Analysis and Testing of Space Structures -- Structural compliance study of an axial-symmetric spacecraft with Vega-C. Static, modal, thermoelastic, buckling, sine test, coupled load, and acoustic analyses in MSC Nastran (ECSS qualification levels).
    \item Aerospace Technologies and Materials -- a) Modeled a composite panel with embedded piezoelectric plate; validated natural frequencies via NX Nastran and testing. b) Optimized an additive manufacturing component with Fusion 360 and Abaqus, achieving 61.63\% volume reduction.
  \end{itemize}}

\vspace{0.6em}

\cventry
  {B.S. in Aerospace Engineering}
  {Politecnico di Milano}
  {Milan, Italy}
  {Sept. 2018 -- Sept. 2021}
  {}

%-------------------------------------------------------------------------------
% SOFTWARE SKILLS
%-------------------------------------------------------------------------------
\section{Software Skills}

\vspace{0.3em}
\noindent\begin{tabularx}{\textwidth}{@{}r X@{}}
  \cvskillrow{CAE}{Simcenter 3D, Abaqus, Femap, MSC Nastran, Patran, Ansys Mechanical, HyperMesh}
  \cvskillrow{CAD}{CATIA, SolidWorks, Solid Edge, Autodesk Inventor}
  \cvskillrow{Programming}{MATLAB, Python, LaTeX}
  \cvskillrow{Office}{Microsoft Office}
\end{tabularx}

%-------------------------------------------------------------------------------
% LANGUAGES
%-------------------------------------------------------------------------------
\section{Languages}

\vspace{0.3em}
\noindent\begin{tabularx}{\textwidth}{@{}r X@{}}
  \cvskillrow{Italian}{Mother tongue}
  \cvskillrow{English}{Proficient user}
  \cvskillrow{French}{Basic user}
\end{tabularx}

\end{document}
